\subsection{INVERSE LIMITS OF COMPACT SPACES}

PROPOSITION 8. Let \((\mathbf{X}_{\alpha}, f_{\alpha \beta})\) be an inverse system of compact spaces indexed by a directed set I such that \(f_{\alpha \alpha}\) is the identity mapping for each \(\alpha \in \mathbf{I}\) . Let \(\mathbf{X} = \varprojlim \mathbf{X}_{\alpha}\) be the inverse limit and \(f_{\alpha}: \mathbf{X} \to \mathbf{X}_{\alpha}\) the canonical mapping (§ 4, no. 4). Then

\begin{enumerate}
\def\labelenumi{\alph{enumi})}
\tightlist
\item
  X is compact and for each \(\alpha \in I\) we have
\end{enumerate}

\[
f _ {\alpha} (\mathbf {X}) = \bigcap_ {\beta \geqslant \alpha} f _ {\alpha \beta} (\mathbf {X} _ {\beta}).\tag{I}
\]

\begin{enumerate}
\def\labelenumi{\alph{enumi})}
\setcounter{enumi}{1}
\tightlist
\item
  If the \(\mathbf{X}_{\alpha}\) are all non-empty then \(\mathbf{X}\) is non-empty.
\end{enumerate}

X is a closed subspace of \(\prod_{\alpha} X_{\alpha}\) (§ 8, no. 2, Proposition 7, Corollary 2)

which is compact by Theorem 3 of no. 5 and Proposition 3 of no. 3. The other assertions are consequences of Set Theory, chap.~III, § 7, no. 4, th. 1. We apply this theorem by taking \(\mathfrak{S}_{\alpha}\) to be the set of closed subsets of \(\mathbf{X}_{\alpha}\) . The conditions (i) and (ii) are just axioms \((\mathrm{O}_{\mathbf{i}}^{\prime})\) and \((\mathbf{C}^{\prime \prime})\) respectively; condition (iii) is satisfied since \(\{x_{\alpha}\}\) is closed and \(f_{\alpha \beta}\) continuous ( \(\S 2\) , no. 1, Theorem 1), and lastly condition (iv) is satisfied by reason of Corollary 2 of Theorem 2 of no. 4.

COROLLARY I. Let \((\mathbf{X}_{\alpha}, f_{\alpha\beta})\) be an inverse system of topological spaces indexed by a directed set, such that for each pair of indices \(\alpha\) , \(\beta\) for which \(\alpha \leqslant \beta\) and for each \(x_{\alpha} \in X_{\alpha}\) , \(\overline{f}_{\alpha\beta}(x_{\alpha})\) is compact. Then equation (1) is valid and \(\overline{f}_{\alpha}(x_{\alpha})\) is compact for each index \(\alpha\) and each \(x_{\alpha} \in X_{\alpha}\) .

For each \(x_{\alpha} \in \bigcap_{\beta \geqslant \alpha} f_{\alpha \beta}(X_{\beta})\) and each \(\beta \geqslant \alpha\) , let \(L_{\beta}\) denote \(\overline{f}_{\alpha \beta}^{1}(x_{\alpha})\) . If \(\alpha \leqslant \beta \leqslant \gamma\) , then we have \(f_{\beta \gamma}(L_{\gamma}) \subset L_{\beta}\) and the set of all indices \(\beta \geqslant \alpha\) is cofinal in the index set. It follows immediately that the \(L_{\beta} (\beta \geqslant \alpha)\) form an inverse system of topological spaces (with the restrictions of the \(f_{\beta \gamma}\) as mappings), whose inverse limit \(L\) is homeomorphic to \(\overline{f}_{\alpha}^{1}(x_{\alpha})\) . Since by hypothesis the \(L_{\beta}\) are compact and not empty, the corollary follows from Proposition 8.

COROLLARY 2. Let \((\mathbf{X}_{\alpha}, f_{\alpha \beta})\) and \((\mathbf{X}_{\alpha}^{\prime}, f_{\alpha \beta}^{\prime})\) be two inverse systems of topological spaces indexed by the same directed set I, and let \((u_{\alpha})\) be an inverse system of mappings \(u_{\alpha}: \mathbf{X}_{\alpha} \to \mathbf{X}_{\alpha}^{\prime}\) . Let \(\mathbf{X} = \varprojlim \mathbf{X}_{\alpha}\) , \(\mathbf{X}^{\prime} = \varprojlim \mathbf{X}_{\alpha}^{\prime}\) , \(u = \varprojlim u_{\alpha}\) . Then:

\begin{enumerate}
\def\labelenumi{\alph{enumi})}
\item
  If \(x' = (x_{\alpha}') \in \mathbf{X}'\) is such that \(\overline{u}_{\alpha}(x_{\alpha}')\) is compact and non-empty for each \(\alpha \in I\) , then \(\overline{u}^1(x')\) is compact and non-empty.
\item
  If the \(X_{\alpha}\) are compact, the \(X_{\alpha}^{\prime}\) Hausdorff and the \(u_{\alpha}\) surjective and continuous, then u is surjective.
\end{enumerate}

Let \(\mathbf{L}_{\alpha}\) denote \(\overline{\boldsymbol{u}}_{\alpha}^{1}(x_{\alpha}^{\prime})\) ; then clearly the \(\mathbf{L}_{\alpha}\) form an inverse system of topological spaces (with the restrictions of the \(f_{\alpha \beta}\) as maps) and \(\overline{\boldsymbol{u}}^{1}(x^{\prime} = \mathbf{L})\) is the inverse limit of the \(\mathbf{L}_{\alpha}\) ; assertion a) therefore follows from Proposition 8. Assertion b) is an immediate consequence, in view of Proposition 3 of no. 3.

